\section{Permanencia cuando quedan al menos \texorpdfstring{$\delta n$}{δn} filas}\label{macro}

\subsection{Proceso de revelado y trayectoria de comparación}

Fijamos $0<\delta<1$. Tras $s$ elecciones escribimos $k=n-s$, $x=k/n$ y $t=s/n$. Bajo $P^{\rm B}$ se
revela la permutación uniforme $\sigma$ de las filas una a una. Condicionalmente a la historia
$\mathcal F_s$ de las elecciones y de las filas ya reveladas, la próxima fila es uniforme entre las
$k$ filas vacías. Si $p$ es el punto de Bethe del tablero actual, la probabilidad condicional de
seleccionar una
casilla $e$ es $p_e/k$. En las casillas fuera del soporte ponemos $p_e=0$. Los índices discretos van de
$1$ a $n$. Los centros físicos
$r_i=(i-\tfrac12)/n$ y $c_j=(j-\tfrac12)/n$ sitúan la diagonal $D(e)=i-j$ en $r_i-c_j$ y la
antidiagonal $A(e)=i+j-1$ en $r_i+c_j$. Denotamos por $X_f$ el indicador de disponibilidad de una
línea y por $Y_e=\prod_{f\ni e}X_f$ el de una casilla. Como en la Sección~\ref{leyes}, $\ell_f$ y
$h_f$ son la masa y la vacancia de una diagonal o antidiagonal $f$ en el punto $p$.

\begin{lema}[Pérdida de una casilla]\label{macro:casilla}
La probabilidad condicional de ocupar una fila o columna disponible es $1/k$. La probabilidad de
ocupar una diagonal disponible $f$ es $\ell_f/k$. Para $e=(i,j)$ se tiene
\[
 \E_{\rm B}[Y_e(s+1)-Y_e(s)\mid\mathcal F_s]
 =-\frac{Y_e(s)}{k}\bigl(2+\ell_{D(e)}+\ell_{A(e)}-3p_e\bigr).
\]
\end{lema}
\begin{proof}
Las dos primeras fórmulas se obtienen sumando $p_e/k$ sobre la línea correspondiente. Si $Y_e=1$, las
cuatro líneas que pasan por $e$ sólo se cortan en esa casilla. En la suma de sus masas $e$ aparece
cuatro veces y hay que contarla una vez. Si $Y_e=0$, ninguna elección puede volver a hacerla
disponible.
\end{proof}

En el cuadrado unidad usamos el perfil
\[
 \nu_1(r,c)=b_0r_R(r)r_C(c)u_D(r-c)u_A(r+c).
\]
Sus marginales de filas y columnas valen uno. Para $d\in[-1,1]$ y $a'\in[0,2]$, sus proyecciones
sobre las dos familias diagonales, es decir, las masas de sus fibras, son
\[
 m_D(d)=\int_{\{0\le r,r-d\le1\}}\nu_1(r,r-d)\,dr,\qquad
 m_A(a')=\int_{\{0\le r,a'-r\le1\}}\nu_1(r,a'-r)\,dr .
\]
Aquí y en adelante la notación bajo el signo de integral exige que ambas coordenadas pertenezcan a
$[0,1]$. Las vacancias son $u_f=1-m_f$. Usamos $c_\nu\le\nu_1\le C_\nu$, $u_f\ge\kappa_\nu>0$ y la
regularidad Lipschitz del perfil lleno \cite[Teoremas~1.1--1.3, Proposición~5.2 y
Teorema~7.7]{perfil}. Las masas de las fibras son Lipschitz en su parámetro. En efecto, desplazarlo en
$\epsilon$ mueve los extremos de integración a lo sumo en $|\epsilon|$ cada uno y cambia el integrando
en la parte común a lo sumo en $L_\nu|\epsilon|$.

La proposición siguiente reúne los objetivos deterministas con los que comparamos el proceso cuando
queda una fracción $x$ de filas. Para una diagonal $f$, $U_f(x)$ es el objetivo de su indicador de
disponibilidad y $\ell_f^*(x)$ el de su masa. En un punto $(r,c)$, $J_x$ es el objetivo del indicador
de disponibilidad de la casilla y $v_x/n$ el del peso de Bethe de una casilla disponible.

\begin{proposicion}[Trayectoria de comparación]\label{macro:trayectoria}
Para una línea diagonal $f$, un punto $(r,c)$ y $0\le x\le1$, escribimos
\begin{gather*}
 U_f(x)=u_f+xm_f,\qquad \ell_f^*(x)=\frac{xm_f}{U_f(x)},\\
 J_x=x^2U_D(x)U_A(x),\qquad
 W_x=\frac{\nu_1}{U_D(x)U_A(x)}.
\end{gather*}
Los subíndices $D,A$ sin argumento se evalúan en $r-c,r+c$, respectivamente. Como $u_f\ge\kappa_\nu$,
$W_x$ se prolonga por continuidad hasta $W_0=\nu_1/(u_Du_A)=b_0r_Rr_C$, el peso usado en la
Sección~\ref{meso:regularidad}. Se tiene $J_1=1$. Para $0<x\le1$, ponemos $v_x=W_x/x$. Entonces
$J_xv_x=x\nu_1$, ambas marginales de $J_xv_x$ valen $x$ y
\begin{equation}\label{macro:derivadas}
 \partial_t U_f=-\frac{\ell_f^*}{x}U_f,\qquad
 \partial_t\log J_x=-\frac{2+\ell_D^*+\ell_A^*}{x},\qquad
 1-\ell_f^*=\frac{u_f}{U_f}\ge\kappa_\nu .
\end{equation}
En particular,
\[
 v_x=\frac{b_0}{x}\,r_Rr_C\frac{u_D}{U_D}\frac{u_A}{U_A},
 \qquad xv_x\le C_\nu/\kappa_\nu^2 .
\]
\end{proposicion}
\begin{proof}
Se tiene $U_f(1)=1$, $\kappa_\nu\le U_f(x)\le1$ y $dx/dt=-1$. Por tanto $\partial_t U_f=-m_f$. Derivar
el producto que define $J_x$ da la segunda identidad de \eqref{macro:derivadas}. La fórmula de $v_x$
se obtiene sustituyendo la factorización de $\nu_1$. La identidad del producto y las afirmaciones
sobre las marginales siguen por integración.
\end{proof}

La comparación que sigue necesita los productos de disponibilidades de varias líneas. Las
observaciones de las dos familias diagonales por separado sólo dan las medias de cada factor.
Definimos a continuación sumas de productos y demostramos su concentración sin imponer independencia a
las líneas.
